
\chapter{The Mathematical Minimum}

Up till now, this book as described quantum computing qualitatively. An introduction to the introduction. To continue this journey, it is essential to give a mathematical description that illustrates the beauty of the symmetry and showcases the power of the theory. To begin with, we must go beyond the limits of reality. 

\section{Real and imaginary numbers}

So far everything in this textbook has been described using real numbers- that is the set numbers you can produce on most calculators. For quantum computing, we require going beyond to complex numbers. 

There are two square roots of $4$: $+2$ \& $-2$. For numbers that aren't integers there are algorithms we can run on our calculators to work out the roots to arbitrary precision. In physics there are many situations in which the only valid solution to an equation calls for the square root of a negative number. To make this work, it is necessary to first define the square root of $-1$ as $i = \sqrt{-1}$. 

Using some basic algebra we can factorise any surd (square root number) into a product of two surds $ \sqrt{ab} = \sqrt{a} \times \sqrt{b}$. This holds true for negative numbers allowing us to take the square root of any negative number by breaking it down into the square root of $-1$ and the positive magnitude of that number. This can be written as:

$$ \sqrt{-a} = \sqrt{-1} \times \sqrt{a} $$

$$ \sqrt{-a} = i \sqrt{a} $$
For example the square root of $-4$ can be written as $+2i$, $-2i$.
These numbers that can be written as $ai$ are known as imaginary numbers. 

\subsection{ Complex numbers}

Imaginary numbers work similarly to the real numbers everyone uses on a  daily basis. One can divide, multiply, add \& subtract them just as you can for real numbers. What this also means is that any of $ \div ,\times, +, - $ operations can be applied to a combination of both real and imaginary numbers. For instance, the positive root of $-4$, $2i$ can be added to $3$ as $3 + 2i $. Any such number that can be written as a sum of a real and imaginary number is called a \textit{complex number}. More generally we can define this mathematically as:

$$ z = a + ib $$

Where $a$ and $b$ are purely real numbers. The real part of $z$ is known as $ Re(z) = a$ and similarly the imaginary component of is $ Im(z) = b$. 

\subsection{ The Argand Diagram}

Complex numbers can be represented on a graph by splitting them into their real \& imaginary parts. In the Argand diagram the x axis represents the real part and the y axis represents the imaginary part. 

If a were negative, the real part would be negative so the arrow would point to the left instead. If b were negative, the imaginary part would be negative, so the arrow would point downwards instead. 

\subsection{ Complex conjugate}

Every complex number has a real \& imaginary part. By taking the imaginary part of the complex number and making it negative, we get the complex conjugate. For example if the complex number $z$ is 

$z = x + iy$ 

The complex conjugate would be $z^*$ written as 

$z^* = x - iy$ 

In the figure below the complex number $z$ is shown along with its complex conjugate. In this diagram $\bar{z}$ is used rather than $z^*$ but it means the same thing.  


\begin{figure}[ht!]
    \centering
    \includegraphics[scale=0.5]{images/Complex_conjugate_picture.svg.png}
    \caption{The complex number $z$ and its complex conjugate $\bar{z}$.  Note how the real part is the same, but the imaginary part is on the negative axis.}
    \label{fig:complex_conjugate}
\end{figure}

 \subsection{ Magnitude of complex numbers }

 It can be very useful to describe how big a complex number is. Since it has a real component and an imaginary component, the magnitude of a complex number indicates how big it is across both these components. 

 The magnitude of a complex number $z$ is denoted by two vertical bars as $|z|$. It is calculated the same was as the Euclidean distance as 

 $$|z| = \sqrt{x^2 + y^2}$$

 With $z = x + yi$. 

 The magnitude of a complex number is important for figuring out the probability of getting an output from a quantum computation. This will be explored in the next chapter.

\subsection{ Interlude on Trigonometry } 

Trigonometry is defined as \textit{the branch of mathematics dealing with the relations of the sides and angles of triangles and with the relevant functions of any angles} but for our use it is much better to consider trigonometry the study of rotations. In later chapters, quantum gates will be described as rotations. Consider a circle with radius 1. Any point on the circle can be reached by rotating on the surface by an angle. The figure below shows how we can reach any point.

\begin{center}
    \includegraphics[scale=0.65]{images/TrigonometryUnitCircle_700.png}
\end{center}

The point P is reached by the angle $\theta$. It has coordinates $x = \cos(\theta)$ and $ y = \sin(\theta)$.

In chapter 5, quantum gates will be described in terms of rotations around a higher dimensional sphere. 

\subsection{ Euler's formula}

\textbf{Euler's formula} is very useful and can be stated as

$$ e^{i \theta} = cos(\theta) + i sin(\theta) $$

Where $\theta$ is an angle in radians. The significance of this is that one can represent any complex number in this fashion simply by multiplying the left-hand side by some magnitude. For example, the complex number $ z = \frac{1}{2}( 1 + \sqrt{3}i)$ 

Can be written as 

$$ z = cos(\pi/3) + i sin(\pi/3) $$

With Euler's formula this becomes 

$$ z = e^{i \pi/3} $$

In the context of quantum information this object $ e^{i\theta}$ is called a phase and is of exceptional importance. To explain why $\theta$ is so significant would require its own chapter. Instead it shall be incorporated into later chapters as needed.


\subsection{ Polar form of complex numbers}

We can represent any complex number on the \textit{Argand diagram} in terms of its real and imaginary parts. From trigonometry, we can also represent any point on a circle using an angle and a given radius. 

Euler's formula allows us to put these two together, representing any complex number as a point on a circle. 

\begin{center}
    \includegraphics{images/Polar_Argand-diagram.png}
\end{center}

There are some complex numbers that are too big to be written as $e^{i\theta}$. For instance, the complex number z = 5 + 12i cannot be written as $cos(\theta)$ because $cos(\theta)$ only goes up to 1. To overcome this, we can multiply by some magnitude to allow for bigger numbers. By multiplying by 13 we can write 

$$ z = 13(cos(1.176) + i sin(1.176) )= 5 + 12i $$

$$ z = 13e^{i 1.176}$$

Where did the 13 come from? 

The factor in front of the $e$ is the magnitude of the complex number. This is the same magnitude described in 3.2.5. This is found by taking the norm of the real and complex part 

$$ 13 = \sqrt{ 5^2 + 12^2} $$

Using Euler's formula, any complex number can be written in terms of its magnitude and an angle. This gives us 

$$z = x + iy = re^{i\theta}$$ 

Where $r$ is the magnitude of the complex number $r = |z| = \sqrt{x^2 + y^2}$ and $\theta$ is the angle made between the real and the imaginary plane

\section{Linear Algebra}

The second part of chapter 3 will cover linear algebra. Linear algebra is the framework for quantum computing, so it is very important to be familiar with. Like trigonometry, linear algebra underpins most of physics, engineering and computer science. Its importance is difficult to overstate. In 1939, Paul Dirac reformulated quantum mechanics using linear algebra [1]. Linear algebra is so important that this way of describing quantum mechanics is known as matrix mechanics. There are two key elements to linear algebra: vectors and matrices. 

\subsection{ Vectors}

Vectors are objects that have both direction \& magnitude. Velocity is a vector because it has a magnitude, the speed at which the object travels with, and a direction, where it's moving towards. We will use a different notation to represent a vector with a straight line and an angled bracket as $ \ket{v} $ .

Imagine a car travelling diagonally across a road. At the same time, the car is travelling along the direction of the road as well as perpendicularly across the road. It can be said that the velocity of the car has a component pointing along the road and another component across the road. What makes vectors useful to work with is the fact that a vector has components. These are the individual magnitudes of the vector in each direction. We can encode the speed of the car along the road ($v_{alon}$)and the speed of the car across the road ($v_{acro}$) in a vector with two components: 

$$
\ket{v} = \begin{bmatrix}v_{alon} \\ v_{acro}
\end{bmatrix}
$$

Here the vector $\ket{v}$ has two components, $v_{alon}$ \& $v_{acro}$ so the vector has a dimension of 2. In 3D space, any point can be described by 3 vectors. In general, a vector can be of any dimension, with any number of components. 

We can write any n-dimensional vector as

$$
\ket{v} = \begin{bmatrix} v_0 \\ v_1 \\ \vdots \\ v_{n-1} \end{bmatrix}
$$

Where the subscript $v_i$ indicates the i $^{th}$ component of $\ket{v}$. 

\textbf{Important note:} We start counting from 0 and end at $n-1$. This convention of counting from 0 is widely used in QC and conveniently is also used in python too. 

\subsection{ Adding vectors }

It can be useful to add vectors. All this requires is adding the corresponding components of each vector. Adding two vectors, $\ket{a}$ \& $\ket{b}$

$$\ket{a} + \ket{b} = \begin{bmatrix} a_0 \\ a_1 \\ \vdots \\ a_{n-1} \end{bmatrix} + \begin{bmatrix} b_0 \\ b_1 \\ \vdots \\ b_{n-1} \end{bmatrix} = \begin{bmatrix} a_0 + b_0 \\ a_1 + b_1 \\ \vdots \\ a_{n-1} + b_{n-1} \end{bmatrix} $$

The same can be done in reverse, a vector can be split into different components. For example 

$$
\ket{a} = \begin{bmatrix} 3 \\ 2 \end{bmatrix} = \begin{bmatrix} 3 \\ 0 \end{bmatrix} + \begin{bmatrix} 0 \\ 2 \end{bmatrix}
$$

This can be useful for expressing a vector in terms of its components. 

\subsection{ Multiplying a vector by a scalar }

Vectors can be made smaller or bigger by enlarging them by some scale factor. To make a vector bigger by some scale factor, the vector is multiplied by a scalar. To scale the vector $\ket{v}$ by some scalar quantity, $a$, this would be done by multiplying each component of the vector by the scalar factor. 

$$ a \ket{v} = a\begin{bmatrix} v_0 \\ v_1 \\ \vdots \\ v_{n-1} \end{bmatrix} = \begin{bmatrix} a \times v_0 \\ a \times v_1 \\ \vdots \\ a \times v_{n-1} \end{bmatrix} $$

\subsection{The inner product}

It can be very useful to compare two vectors by measuring how aligned they are. Imagine two vectors pointing the same way. These two vectors would have a lot in common, and so would be perfectly aligned. The measure of how similar two vectors are is the inner product. The inner product can be calculated by multiplying each component of the vectors and adding them up. 

For example take the vectors $\ket{a}$ \& $\ket{b}$ as 

 $$
 \ket{a} = \begin{bmatrix} 1 \\ 2 \end{bmatrix} ;  \ket{b} = \begin{bmatrix} 3 \\ 1 \end{bmatrix} 
 $$

 The inner product between them is $\braket{a|b}$. This can be computed by taking the column vector $\ket{a}$ and turning it into a row vector where every element of $\ket{a}$ has taken the complex conjugate. 

 $$
\braket{a|b} = \begin{bmatrix} 1^* & 2^* \end{bmatrix}  \begin{bmatrix} 3 \\ 1 \end{bmatrix}
 $$

Here the 1st component of $\bra{u}$ is multiplied by the 1st component of $\ket{v}$ and the second component of of $\bra{u}$ is multiplied by the second component of $\ket{v}$.

$$
= (1 \times 3) + (2 \times 1) = 3 + 2 = 5 
$$

This gives a scalar quantity from two vectors and is used to indicate how aligned they are. The figure below shows this, where the length of $\ket{a}$ shows the component of $\ket{u}$ in the direction of $\ket{v}$. This is then scaled by the length of $\ket{v}$ to give the inner product . 

\includegraphics[scale=0.5]{images/Dot_product_visualisation.png}


The inner product depends on the angle between the vectors. If the vectors are pointing the same way, it will be bigger, if the vectors are pointing in completely opposite directions, it will be smaller. 

\subsection{ Perpendicular vectors }

If two vectors have no components in common, their inner product will be zero. This can be shown by two vectors that have an angle of $ 90^{\circ} $ between them.  We call these vectors perpendicular. In quantum mechanics, the preferred term is \textbf{orthogonal} which means the same thing but applies to more than just vectors. 

For example the vectors $\ket{p} = \begin{bmatrix} 1 \\ 0 \end{bmatrix}$ ; $\ket{q} = \begin{bmatrix} 0 \\ 1 \end{bmatrix}$ have a inner product  of 0. 

$$ \begin{bmatrix} 1^* & 0^* \end{bmatrix}  \begin{bmatrix} 0 \\ 1 \end{bmatrix} = (1 \times 0) + (0 \times 1) = 0$$

\section{ Matrices}

Another object, closely related to a vector, is a matrix. A matrix is just an array of numbers which you can perform a matrix product on. All the quantum gates will be represented as matrices, so understanding how they work will be essential for quantum computations. Matrices are usually denoted by capital letters, for example $M$

$$ M = \begin{bmatrix} 1 & 2  \\ 3 & 4 \end{bmatrix} $$

$M$ has 4 elements: 1,2,3 \& 4. 

Each element in $M$ has an index: a pair of numbers that gives the position in the matrix. This is a bit like how map coordinates indicate a position. Starting from the top left, the index increases from left to right or from top to bottom. For example, the bottom left element (3) of $M$ would have the index $M_{10}$ as the first row, zeroth column. The choice of starting from zero is made to be consistent with binary and is also useful for programming languages like python that are widely used in quantum computing. 

$M$ could also be written as a collection of indices

$$ M = \begin{bmatrix} M_{00} & M_{01}  \\ M_{10} & M_{11} \end{bmatrix} $$


\subsection{ Matrix-vector product} 


Matrices, or more generally 2D arrays of numbers, are useful for storing information. The pixels on a screen have colour intensity values that can be efficiently stored as a matrix where the number of rows and columns is equal to the horizontal and vertical pixels on a screen (1080 by 1920 is a popular choice). As well as storing information, matrices can be used to transform vectors. This is done using the matrix-vector product. The matrix $M$ acting on the vector $\ket{a}$ is denoted as $M\ket{a}$. The matrix multiplication is done by multiplying the corresponding elements of $M$ \& $\ket{a}$. Starting from the first row of $M$ we multiply out the elements of $M$ with the elements of $\ket{a}$ starting from the left and add them up. Let's call the transformed vector $\ket{a'}$. We can compute $\ket{a'}$ by doing the matrix-vector product of $M$ on $\ket{a}$ as

$$ \ket{a'} = M\ket{a} $$

Explicitly calculating this gives us

$$ \ket{a'} =  \begin{bmatrix} 1 & 2  \\ 3 & 4 \end{bmatrix} \begin{bmatrix} 1 \\ 2 \end{bmatrix} $$

$$ = \begin{bmatrix} (1 \times 1) + (2 \times 2) \\ (3 \times 1) + (4 \times 2) \end{bmatrix} $$

$$ \ket{a'} = \begin{bmatrix} 5 \\ 11 \end{bmatrix} $$
Matrix multiplication follows some of the same relationships as multiplying numbers, but not all of them. 

\subsection{Multiple matrices}

A matrix transforming a vector can be thought of as an instruction. The original vector represents some initial starting position (like a coordinate x,y) and the matrix maps that point onto a new point (x', y'). Matrices can be applied in a sequence. Starting from the vector $\ket{a}$, applying the matrix $A$ followed by $B$ would look like 

$$ BA\ket{a} $$

In linear algebra the order is read from right to left, whatever is the furthest to the right happens first and whatever is furthest to the left happens last. Matrix multiplication follows some of the same relationships as multiplying numbers, but not all of them. In regular multiplication, the order in which numbers are multiplied doesn't matter. $ 5 \times 3 = 3 \times 5 = 15 $. For any numbers $a, b$, $a \times b = b \times a$. The same cannot be said for any two matrices, $ AB \neq BA $. This relationship, whether the order of operations matters, is defined as whether $A$ and $B$ commute. Integer multiplication is commutative, matrix multiplication does not usually commute.  

\subsection{ Inverse matrices }

For many matrix-vector products, we can go back and get the original vector from the transformed vector. This means, if we know the transformed vector $\ket{a'}$ \&  the matrix $M$ we can work out the original vector $\ket{a}$. 

To get the original vector, we do the matrix-vector product, but with the transformed vector and the inverse of the matrix. This gives us 

$$
M^{-1}\ket{a'} = \ket{a}
$$

By using our equation for $\ket{v'} = M\ket{a}$ from earlier, this gives us 

$$ M^{-1}M\ket{a} = \ket{a} $$

Notice how $\ket{a}$ is completely unchanged from applying $M$ and then $M^{-1}$. The operation $M^{-1}M$ effectively does nothing. The same would be true if we applied $M^{-1}$ first and then $M$ as $MM^{-1}$. The general name for the "do nothing" matrix is the identity. The identity is

$$M^{-1}M = MM^{-1} = I $$

Before we had $\ket{a'} = M\ket{a}$. And now we would like to go backwards from $\ket{a'}$ to $\ket{a}$. This is done using the inverse of $M$: $M^{-1}$

$$ M^{-1}\ket{a'} = \ket{a} $$

Notice how we can take the original equation and do the matrix multiplication by $M^{-1}$ and get the same result

$$ M^{-1}\ket{a'} = M^{-1}M\ket{a} = \ket{a} $$

$M^{-1}$ is cancelling out $M$ so that there is no effect on $\ket{a}$. The matrix that represents this "no effect" is called the identity $I$ defined by 

$$ I\ket{v} = \ket{v}$$

As a matrix, the identity looks like 

$$ I = \begin{bmatrix} 1 & 0 \\ 0 & 1 \end{bmatrix}$$

\subsection{ Unitary matrices} 

For all quantum gates, an important property they have is that they are unitary. Unitary matrices have an inverse equal to their conjugate transpose (the dagger at the top).

$$ U^\dagger = U^{-1} $$

This means 

$$ U^\dagger U = U^{-1}U = I $$

The conjugate transpose is the same process as it is to turn a ket (column vector) into a bra (row vector). Take You swap the rows and the columns and then replace every imaginary number with the negative version. 

All quantum gates are unitary matrices, it is very common to see the term unitary when talking about any quantum circuit. 

\section{Chapter Exercises}
\begin{enumerate}
    \item Consider the complex numbers 

    $$ z = 4 + 5i ; y = -3 + i; m = -2 - 2i $$
    \begin{enumerate}[i]
        \item What is $Re(z)$? 
        \item What is $Im(m)$? 
        \item Sketch z,y,m on the Argand diagram
    \end{enumerate}

    \item Consider the complex number $z = 3 + 5i$ . 
        \begin{enumerate}[i]
            \item What is the complex conjugate of $z$, $z^*$?
            \item What is the magnitude of $z$?
            \item What is the magnitude of $z^*$?
        \end{enumerate}

    \item Compute the inner product for the vectors 

        $$ \ket{w} = \frac{1}{\sqrt{2}} \begin{bmatrix} i \\ 1 \end{bmatrix} ; \ket{z} = \frac{1}{\sqrt{3}} \begin{bmatrix} \sqrt{2} \\ i \end{bmatrix}$$

        i. $\braket{w|z}$

        ii. $\braket{z|w}$

        iii. Bonus question: Compare i. ii. . Is the result the same as from Exercise 3.2 iv? 
        If not, what has changed?


    \item Imagine the vector $\ket{v}$ given by $$\ket{v} = \begin{bmatrix} 2 \\ 1 \end{bmatrix}$$
        \begin{enumerate}[i]
            \item What is the $0^{th}$ component of $\ket{v}$ ? What about the $1^{st}$ component?
            \item Write $\ket{v}$ in the form $$ \ket{v} = a\begin{bmatrix} 1 \\ 0 \end{bmatrix} + b \begin{bmatrix} 0 \\ 1 \end{bmatrix} $$
            \item Write down the vector $5\ket{v}$ and describe what has changed to the vector  
            
        \end{enumerate}

    \item Let the vector $\ket{q}$ be given as 

    $$ \ket{q} = \frac{1}{4} \begin{bmatrix} -3 \\ 4 \end{bmatrix} $$

    \begin{enumerate}[i]
        \item Using $\ket{v}$ from the last exercise, compute  $\ket{y} = \ket{v} + \ket{q} $
        \item Compute the inner product $ \braket{q | v} $
        \item Compute the inner product $ \braket{v | q} $ 
        \item Bonus question: Compare your answers for ii. \& iii. Explain any differences, if any, or why there is no difference. Does this hold for all pairs of vectors?
        \item Using the result from either ii. or iii. are the vectors $\ket{v}$ \& $\ket{q}$ perpendicular?
    
    \end{enumerate}
\end{enumerate}

\section{Chapter Summary }

\begin{itemize} 
    \item Vectors are a collection of numbers stored in a row or column
    \item Vectors can be added to each other or multiplied by a scalar
    \item The inner product of two vectors tells us how aligned they are
    \item  Perpendicular (orthogonal) vectors have an inner product of 0
    \item  Matrices are arrays of numbers that can act on vectors 
    \item  Applying a matrix to a vector turns the vector into another vector
    \item  Matrices can have an inverse which does the reverse of the matrix 
    \item  Unitary matrices have an inverse equal to their conjugate transpose
\end{itemize}